\chapter{Introduction}
\label{chapter:introduction}

\section{Project Background}

Phishing is a social-engineering attack in which an attacker imitates a trusted
service to persuade a user to reveal information such as a password, bank-card
number, or account code. A phishing attempt may be delivered through email,
text messages, social media, advertisements, or a compromised website. In many
cases, the user is directed to a web address that looks believable but leads to
a fraudulent page.

The URL is useful for early detection because it can be examined before a user
submits information or downloads content. URL-based detection is also suitable
for a browser extension because it does not require the complete webpage to be
downloaded and analysed. Traditional blacklist systems remain useful for known
threats, but a newly created phishing URL may not yet appear in a blacklist.
Machine-learning models can instead learn structural and character patterns
that separate legitimate and phishing URLs \citep{sahingoz2019machine,
kamran2022conditional}.

The main difficulty is that phishing URLs change over time. Attackers can alter
domain names, subdomains, paths, query strings, spelling, special characters,
and hosting services. At the same time, legitimate browsing also contains long
and complex URLs. A detector that treats every unfamiliar or complicated URL as
phishing may achieve high phishing recall while producing too many false
warnings for normal users.

Generative Adversarial Networks offer a possible way to create additional
training examples. A GAN contains a generator that produces synthetic samples
and a discriminator that learns to distinguish generated samples from real
ones. GAN-based phishing research has considered both generation and detection,
including conditional generation, black-box detector evasion, and comparisons
of several GAN families \citep{kamran2022conditional,bac2021pwdgan,
pham2023evaluation}. This project studies whether generated phishing URLs can
be used defensively as controlled training augmentation.

\section{Problem Statement}

Supervised phishing detectors normally learn from labelled historical data.
Their reported accuracy can be high when the test data closely resembles the
training data. However, performance may fall when URLs are collected from a
different source, when legitimate browsing contains structures that were not
represented in training, or when phishing campaigns adopt new combinations of
known patterns.

Simply adding more generated phishing URLs does not guarantee improvement.
Synthetic samples may be invalid, repetitive, unrealistic, too similar to the
training data, or limited to a small number of templates. They may also change
the class balance in a way that increases missed phishing URLs or false warnings.
The effect of augmentation must therefore be tested under controlled conditions
where the real training data, validation data, test data, classifier settings,
and evaluation metrics remain fixed.

The project addresses the following problem:

\begin{quote}
\emph{How can realistic and varied GAN-generated phishing URLs be used to
improve URL classification without hiding increases in false positives, false
negatives, or failures on URLs collected from different sources?}
\end{quote}

The problem has both a research part and an implementation part. The research
part compares real-only training with GAN-augmented training. The implementation
part places the selected model in a browser extension and studies how model
scores should be shown to users without presenting every uncertain URL as
confirmed phishing.

\section{Aim}

The aim of this project is to design, implement, and evaluate a phishing URL
detection system that uses GAN-generated phishing URLs as controlled data
augmentation, while maintaining a practical balance between phishing detection
and false warnings on legitimate URLs.

\section{Objectives}

The objectives originate from the approved project proposal and are organised
into the measurable tasks shown in Table~\ref{tab:introduction-objectives}.

\begin{table}[htbp]
\centering
\caption{Project objectives and how each objective is evaluated.}
\label{tab:introduction-objectives}
\small
\begin{tabular}{@{}p{0.07\textwidth}p{0.47\textwidth}p{0.36\textwidth}@{}}
\toprule
No. & Objective & Evaluation evidence \\
\midrule
O1 & Prepare labelled legitimate and phishing URL data with duplicate and overlap checks. & Dataset counts, class balance, URL-length analysis, and split checks. \\
O2 & Train traditional and deep-learning real-only classifiers. & RF, SVM, CNN, RNN, LSTM, and hybrid model metrics. \\
O3 & Train GAN models using phishing URLs from the training data only. & Generator and discriminator training records. \\
O4 & Measure generated URL quality and diversity. & Validity, uniqueness, overlap, length, structure, and classifier-challenge tests. \\
O5 & Compare real-only and GAN-augmented training fairly. & Fixed augmentation percentages, identical real rows, fixed test sets, and repeated seeds. \\
O6 & Test generalisation on URLs not used for training. & Internal, external, low-phishing, long-URL, and generator-challenge evaluations. \\
O7 & Select and implement a model for real-time use. & ONNX export, browser inference checks, warning policy, and user testing. \\
\bottomrule
\end{tabular}
\end{table}

\section{Research Questions}

The work is guided by the following research questions:

\begin{description}
\item[\textbf{RQ1:}] How well do traditional, neural, and hybrid classifiers distinguish
legitimate and phishing URLs when trained and tested on fixed real data?

\item[\textbf{RQ2:}] Can the selected GANs generate valid, unique, varied, and believable
phishing-style URLs without copying the real training set?

\item[\textbf{RQ3:}] Does adding fixed percentages of generated phishing URLs improve
accuracy, precision, recall, F1-score, ROC-AUC, FPR, or FNR compared with the
matching real-only model?

\item[\textbf{RQ4:}] Do any improvements remain when the models are evaluated on URLs from
different sources and on test sets where legitimate URLs greatly outnumber
phishing URLs?

\item[\textbf{RQ5:}] Can the selected model run locally in a browser extension while keeping
URL processing consistent with the trained model and limiting unnecessary
navigation blocks?
\end{description}

\section{Project Approach}

The project follows the stages in Figure~\ref{fig:introduction-workflow}. The
stages are kept separate so that generated data cannot enter validation or test
sets and test results cannot be used to change the training procedure.

\begin{figure}[htbp]
\centering
\setlength{\fboxsep}{6pt}
\newcommand{\workflowbox}[1]{\fbox{\parbox{0.78\textwidth}{\centering\small #1}}}
\workflowbox{Collect, clean, label, and split real URLs}\\[-0.1em]
$\downarrow$\\[-0.1em]
\workflowbox{Train real-only baseline and hybrid classifiers}\\[-0.1em]
$\downarrow$\\[-0.1em]
\workflowbox{Train URL GANs and check generated URL quality}\\[-0.1em]
$\downarrow$\\[-0.1em]
\workflowbox{Create fixed GAN-augmented training conditions}\\[-0.1em]
$\downarrow$\\[-0.1em]
\workflowbox{Evaluate real-only and augmented models on unchanged test sets}\\[-0.1em]
$\downarrow$\\[-0.1em]
\workflowbox{Export the selected model and test the browser extension}
\caption{High-level workflow of the project.}
\label{fig:introduction-workflow}
\end{figure}

The main classifier input is the URL text itself. Neural models learn from
characters, while traditional models use lexical measurements such as URL
length, digit ratio, special-character ratio, subdomain count, path depth, and
query structure. The final hybrid model combines a character CNN with a fixed
lexical feature vector. Models are trained with repeated random seeds where
appropriate, and thresholds are selected using development or calibration data
rather than final test labels.

The GAN stage uses phishing URLs from the training data only. Generated
candidates must pass validity, duplication, overlap, structure, and diversity
checks before they can be added to a classifier training set. Augmented models
are then compared with real-only controls using accuracy, precision, recall,
F1-score, ROC-AUC, FPR, and FNR.

\section{Project Scope}

The scope is limited to URL-based phishing detection. This keeps the experiments
reproducible and allows the selected model to operate locally in a browser.

\subsection{Included in the Project}

\begin{itemize}
\item collection and cleaning of labelled legitimate and phishing URLs;
\item traditional, neural, and hybrid URL classifiers;
\item character-level and lexical URL features;
\item GAN training and generated URL quality evaluation;
\item controlled real-only and GAN-augmented comparisons;
\item internal and external evaluation with complete classification metrics;
\item threshold analysis for false-positive and false-negative trade-offs; and
\item a local browser extension with warning and passive review behaviour.
\end{itemize}

\subsection{Outside the Project Scope}

\begin{itemize}
\item analysis of email text, attachments, or sender behaviour;
\item webpage screenshots, visual similarity, or full Document Object Model analysis;
\item automatic submission of visited URLs to a remote monitoring service;
\item a claim that generated URLs predict unknown future attack methods;
\item commercial deployment or certification as a complete security product; and
\item replacement of browser reputation services, DNS security, or human judgement.
\end{itemize}

\section{Main Contributions}

The project makes the following contributions:

\begin{itemize}
\item a controlled comparison of traditional, neural, and hybrid phishing URL
classifiers using clearly separated training and evaluation data;

\item an evaluation of several GAN approaches that separates generated URL
quality from their effect on classifier performance;

\item fixed-percentage augmentation experiments that preserve the real training
rows and use the same evaluation sets for real-only and augmented conditions;

\item an analysis of distribution shift, showing why strong results on an
internal test may not transfer to URLs collected from other sources;

\item a hybrid model that combines character learning and lexical URL features
with multi-source training and score calibration; and

\item a working browser extension that performs local inference and separates
model warnings from automatic navigation blocking.
\end{itemize}

\section{Report Structure}

The complete report is planned as follows:

\begin{itemize}
\item \textbf{Chapter 1 -- Introduction} presents the problem, aim, objectives,
research questions, scope, approach, and contributions.

\item \textbf{Chapter 2 -- Background} explains phishing URLs, URL structure,
classification models, GANs, augmentation, distribution shift, and evaluation
metrics.

\item \textbf{Chapter 3 -- Related Work} reviews recent phishing URL detection
and GAN-based URL generation research and identifies the gap addressed by this
project.

\item \textbf{Chapter 4 -- Methodology and System Design} describes the data,
splits, models, GAN designs, quality checks, augmentation rules, and evaluation
method.

\item \textbf{Chapter 5 -- Implementation} explains the important software
components, training pipeline, GPU use, stored model files, and reproducibility
controls.

\item \textbf{Chapter 6 -- Experiments and Results} presents the baseline,
GAN-quality, augmentation, external, and threshold results with their analysis.

\item \textbf{Chapter 7 -- Browser Extension and User Testing} describes the
browser implementation, warning modes, local data handling, technical tests,
and user feedback.

\item \textbf{Chapter 8 -- Conclusion} answers the research questions and
summarises the main findings and limitations.

\item \textbf{Chapter 9 -- Future Work} presents practical improvements and
research that could continue after this project.
\end{itemize}
